\section[Holomorphic automorphisms of the Deligne--Hitchin moduli space]{Holomorphic automorphisms of the Deligne--Hitchin\\ moduli space}

\subsection{Connected component of the holomorphic automorphism group}\label{aut0mdh}

The smooth locus of the moduli space of f\/lat connections on a compact Riemann surface is equipped with a natural hyper-K\"ahler structure; see~\cite{H} for the case of ${\rm SL}(2,\mathbb C)$-connections, and \cite{GX} for a detailed treatment of the case of f\/lat line bundles.

A hyper-K\"ahler manifold is a Riemannian manifold $(M, g)$ which is K\"ahler for three complex structures $I$, $J$ and $K$ satisfying the quaternionic relations
\begin{gather*}IJ = -JI = K .\end{gather*} The interplay between the Riemannian, symplectic and complex geometric aspects of hyper-K\"ahler manifolds makes them fascinating mathematical objects. A hyper-K\"ahler mani\-fold~$M$ admits a twistor space which contains all of the information about $M$ in a~complex geometric fashion: The compatible complex structures on $M$ are parametrized by $\mathbb CP^1 = S^2 = \big\{(x_1, x_2, x_3) \in {\mathbb R}^3 \,|\, x^2_1+x^2_2+x^2_3 = 1\big\}$ via
\begin{gather*}(x_1, x_2, x_3) \longmapsto I_{(x,_1,x_2,x_3)} = x_1 I+x_2 J+x_3 K ,\end{gather*}
where $I$, $J$, $K$ are the three original complex structures on~$M$. The twistor space $\mathcal Z$ is a~complex manifold with a holomorphic surjective submersion to $\mathbb CP^1$ such that the f\/iber over every \smash{$z \in {\mathbb C}P^1$} is identif\/ied with~$(M, I_z)$. The twistor lines are the ``constant'' sections of the above f\/ibration ${\mathcal Z} = M\times\mathbb CP^1 \longrightarrow {\mathbb C} P^1$. These sections are holomorphic and the normal bundle of a~twistor line is isomorphic to ${\mathcal O}_{\mathbb CP^1}(1)^{\oplus d} \longrightarrow \mathbb CP^1$, where $d = \frac{1}{2}\dim_{\mathbb R}M$. This normal bundle is equipped with additional structures which provide all of the information needed to reconstruct the hyper-K\"ahler structure on~$M$; see~\cite{HKLR} for details.

For a complex reductive group $G_{\mathbb C}$, the hyper-K\"ahler structure on the moduli space of f\/lat $G_{\mathbb C}$-connections on a compact Riemann surface $X$ is given by non-abelian Hodge theory. The complex structure $J$ is induced by the complex Lie group $G_{\mathbb C}$ through the identif\/ication of the moduli space with the Betti moduli space $\operatorname{Hom}(\pi_1(X), G_{\mathbb C})/\!\!/G_{\mathbb C}$. Via the solutions of Hitchin's self-duality equation, the moduli space of stable $G_{\mathbb C}$-Higgs bundles, i.e., the Dolbeault moduli space, gets identif\/ied with the moduli space of irreducible f\/lat connections, yielding the complex structure~$I$. The reverse map ${\mathcal M}_{dR}  \longrightarrow \mathcal M_{\rm Dol}$ arises from Donaldson's twisted harmonic maps construction \cite{Do} for the case of $G_{\mathbb C}={\rm SL }(2,\mathbb C)$, see~\cite{Co} for the general case. It was f\/irst shown by Hitchin in the case of ${\rm SL}(2,\mathbb C)$ that~$I$,~$J$ and $K := IJ$ give rise to a~hyper-K\"ahler structure for the natural Riemannian metric. In the abelian case of $G_{\mathbb C} = \mathbb C^*$ the theory simplif\/ies considerably boiling down to the classical abelian Hodge theory~\cite{GX}.

It was f\/irst noticed by Deligne that the twistor space of the moduli space of f\/lat connections on a Riemann surface $X$ has a convenient description by gluing the Hodge moduli spaces for~$X$ and~$\overline X$ via the Riemann--Hilbert isomorphism; see~\cite{Si95,Si08} for details. We recall below this construction for the rank one case (the case of our interests).

Take a point $(\lambda, L, D) \in \mathcal M_{\rm Hod}(X) = \mathcal M^X_{\rm Hod}(1)$ with $\lambda \neq 0$ and consider the f\/lat connec\-tion~$D/\lambda$ on $L$. Let
\begin{gather*}\rho \colon \ \pi_1(X)  \longrightarrow H^1(X, {\mathbb Z}) \longrightarrow \mathbb C^*\end{gather*}
be the monodromy representation for~$D/\lambda$. After identifying $H^1(X, {\mathbb Z})$ with $H^1(\overline{X}, {\mathbb Z})$ by the identity map of the underlying $C^\infty$ manifolds, the homomorphism~$\rho$ gives a~holomorphic rank one $1$-connection $(1, E, \widetilde D)$ on~$\overline{X}$. Consider the morphism~$f$ in~\eqref{lf}. Let
\begin{gather*}
f_{\overline X} \colon \ {\mathcal M}_{\rm Hod}(\overline{X}) := {\mathcal M}^{\overline X}_{\rm Hod}(1)  \longrightarrow {\mathbb C}
\end{gather*}
be the morphism obtained by substituting $\overline X$ in place of~$X$ in \eqref{lf}. Def\/ine a holomorphic map
\begin{gather}\label{hodge-gluing}
\varphi \colon \  f^{-1}({\mathbb C}^*)  \longrightarrow f^{-1}_{\overline X}({\mathbb C}^*) ,\qquad (\lambda, L, D) \longmapsto \big(\lambda^{-1}, E, \widetilde{D}/\lambda\big)
\end{gather}
of the restrictions of the f\/iber bundles ${\mathcal M}_{\rm Hod}(X)$ and ${\mathcal M}_{\rm Hod}(\overline{X})$ to ${\mathbb C}^* \subset \mathbb C$. This $\varphi$ is clearly a~biholomorphism.

\begin{Definition}\label{de1} The rank one Deligne--Hitchin moduli space \begin{gather*}{\mathcal M}_{\rm DH} = {\mathcal M}_{\rm DH}(X)\end{gather*} is
\begin{gather*}{\mathcal M}_{\rm DH}(X) = \mathcal M_{\rm Hod}(X)\cup_\varphi\mathcal M_{\rm Hod}(\overline{X}).\end{gather*}
\end{Definition}

We note that the map $f\vert_{f^{-1}({\mathbb C}^*)}$ coincides with the composition
\begin{gather*}
f^{-1}({\mathbb C}^*) \stackrel{\varphi}{\longrightarrow} f^{-1}_{\overline X}({\mathbb C}^*) \stackrel{f_{\overline X}}{\longrightarrow}
 {\mathbb C}^* \stackrel{z\mapsto z^{-1}}{\longrightarrow} {\mathbb C}^* .
\end{gather*}
Therefore, $f$ and the composition
\begin{gather*}
\mathcal M_{\rm Hod}(\overline{X}) \stackrel{f_{\overline X}}{\longrightarrow}  {\mathbb C} \stackrel{z\mapsto z^{-1}}{\longrightarrow}
{\mathbb C}P^1\setminus\{0\}
\end{gather*}
patch together to produce a morphism
\begin{gather}\label{mapf}
f \colon \ {\mathcal M}_{\rm DH} \longrightarrow \mathbb CP^1 .
\end{gather}
For any $\lambda \in \mathbb CP^1$, the f\/iber $f^{-1}(\lambda)$ will be denoted by $f^\lambda$. The f\/iber $f^{\infty}$ is the moduli space of Higgs line bundles of degree zero on~$\overline X$. This $f$ is the twistor f\/ibration mentioned earlier.

\begin{Definition} The {\it degree} of a smooth map $i \colon \Sigma \longrightarrow {\mathcal M}_{\rm DH}(X)$ from a compact oriented surface $\Sigma$ is the degree of the composition $f\circ i \colon \Sigma \longrightarrow {\mathbb C}P^1$, where $f$ is the projection in~\eqref{mapf}.
\end{Definition}

\begin{Lemma}\label{lemdeg} Let $\Sigma$ be a compact Riemann surface of genus $g_\Sigma$ and $i \colon \Sigma \longrightarrow {\mathcal M}_{\rm DH}(X)$ be a~holomorphic immersion. Then the degree of the corresponding holomorphic normal bundle \begin{gather*}{\mathcal N} = (i^*T{\mathcal
M}_{\rm DH})/T\Sigma\end{gather*} is
\begin{gather*}
\deg ({\mathcal N}) = (2 g_X +2)\deg (i)-\deg (T\Sigma) = (2 g_X +2)\deg (i) +2g_\Sigma-2 ,\end{gather*}
where $g_X$ is the genus of $X$.
\end{Lemma}

\begin{proof} As in the computation in \cite[pp.~555--556]{HKLR} it can be shown that the tangent bundle of the complex manifold ${\mathcal M}_{\rm DH}$ is holomorphically isomorphic to the pull-back
\begin{gather*}f^*\big({\mathcal O}_{{\mathbb C}P^1}(2)\oplus \big({\mathcal O}_{{\mathbb C}P^1}(1)^{\oplus 2g_X}\big)\big) ,\end{gather*}
where the f\/irst summand is the tangent bundle of $\mathbb CP^1$ and the direct sum ${\mathcal O}_{{\mathbb C}P^1}(1)^{\oplus 2g_X}$ corresponds to the hyper-K\"ahler structure on $\mathcal M_{dR}$. Now the lemma follows immediately from the fact that $\text{degree}(i^*f^*{\mathcal O}_{{\mathbb
C}P^1}(1)) = \text{degree}(i)$.
\end{proof}

We will compute $\operatorname{Aut}(\mathcal M_{\rm DH})_0$ by proving a series of lemmas.
\begin{Lemma}\label{sections}
Let $\alpha, \beta, \omega, \eta \in H^0(X, K_X)$ be holomorphic $1$-forms. There exists a holomorphic section $s \colon {\mathbb C}P^1 \longrightarrow {\mathcal M}_{\rm DH}$ defined by
\begin{gather*}s(\lambda) = [\lambda, \overline{\partial}+\overline{\omega}+\lambda\overline{\eta}, \lambda(\partial+\alpha)+\beta] \in \mathcal M_{\rm Hod}(X)
 \subset {\mathcal M}_{\rm DH}\end{gather*}
for $\lambda \in {\mathbb C} \subset \mathbb CP^1$. Conversely, every section of $f \colon {\mathcal M}_{\rm DH} \longrightarrow {\mathbb C}P^1 $ is of this form.
\end{Lemma}

\begin{proof}
The f\/irst part of the lemma is easily verif\/ied. In order to prove the converse direction we start with a lift of the family of holomorphic structures on ${\mathbb C} \subset {\mathbb C}P^1$: As $\mathbb C$ is simply connected, there exists a~map
\begin{gather*}
\lambda \longmapsto \overline{\omega}_1(\lambda) \in \overline{H^0(X, K_X)} , \qquad \lambda \in \mathbb C\end{gather*}
such that \begin{gather*}\phi\circ s(\lambda) = [\overline{\partial}+\overline{\omega}_1(\lambda)]\end{gather*}
for all $\lambda \in {\mathbb C}$, where~$\phi$ is the projection in \eqref{phi} and $[-]$ denotes gauge equivalence class. Let $\beta \in H^0(X, K_X)$ be the (well-def\/ined) Higgs f\/ield of $s(0)$. A lift of $s$ is then given by
\begin{gather*}\widehat{s} (\lambda) = (\lambda,\overline{\partial}+\overline{\omega}_1(\lambda), \beta+\lambda(\partial+\alpha_1(\lambda))\end{gather*}
for some holomorphic map \begin{gather*}\alpha_1 \colon \ \mathbb C \longrightarrow \Omega^{(1,0)}(X) .\end{gather*} The image of $\alpha_1$ is contained in $H^0(X, K_X) \subset \Omega^{(1,0)}(X)$ because the $\lambda$-connections are integrable. Thus, $\lambda \longmapsto s(\lambda)$, $\lambda \in \mathbb C^*$, produces the following family of f\/lat connections on~$X$:
\begin{gather*}\widehat{\nabla}^\lambda = d+\lambda^{-1}\beta+\overline{\omega}_1(\lambda)+ \alpha_1(\lambda) .\end{gather*}

We can repeat this over $\mathbb CP^1\setminus\{0\}$ for the Riemann surface $\overline X$. Indeed, there exists a lift $\widetilde s$ of the form
\begin{gather*}\widetilde{s}(\mu) = (\mu,\partial +\alpha_2(\mu), \mu(\overline{\partial}+ \overline{\omega}_2(\mu))+\overline{\eta})\end{gather*}
with $\mu = \frac{1}{\lambda}$. The corresponding f\/lat connections are denoted by $\widetilde{\nabla}^\lambda$. By the gluing condition in Def\/inition \ref{de1} there exists a~gauge transformation $g(\lambda) \colon X \longrightarrow \mathbb C^*$ for every $\lambda \in \mathbb C^*$ such that \begin{gather*}\widehat{\nabla}^\lambda\cdot g(\lambda) = \widetilde{\nabla}^\lambda ;\end{gather*}
this gauge transformation is unique up to a constant scalar. Since the connection 1-forms of $\widehat{\nabla}^\lambda$ and $\widetilde{\nabla}^\lambda$ are both harmonic 1-forms, their dif\/ference
\begin{gather*}\widehat{\nabla}^\lambda-\widetilde{\nabla}^\lambda = \chi(\lambda)\end{gather*}
is a lattice point
\begin{gather}\label{Gamma}
\chi(\lambda) \in \Lambda := \left\{\!\chi \in \Omega^1(X, \mathbb C) \,|\, d\chi = d*\chi = 0;
\int_\gamma\chi\in2\pi\sqrt{-1}\mathbb Z \text{ for all closed curves } \gamma\!\right\}\!.\!\!\!\!\end{gather}
As the connection 1-forms depend continuously on~$\lambda$, it follows that $\chi$ is independent of $\lambda$. By construction, the connection 1-form of $\lambda \longmapsto \widetilde{\nabla}^\lambda$ has a f\/irst order pole at $\lambda = \infty$ whose residue is~$\overline{\eta}$.
This implies that $\overline{\omega}_1(\lambda)$ must be a polynomial of degree at most one and that $\alpha_1(\lambda)$ is a~constant~$\alpha$.
\end{proof}

\begin{Lemma}\label{tau-lemma} Let $T \colon {\mathcal M}_{\rm DH} \longrightarrow {\mathcal M}_{\rm DH}$ be a~holomorphic automorphism. Then $T$ maps fibers of~$f$ to fibers of~$f$. Moreover, it covers the automorphism $\lambda \longmapsto \tau\lambda$ or $\lambda \longmapsto \frac{\tau}{\lambda}$ of ${\mathbb C}P^1$ for some $\tau \in \mathbb C^*$; if the latter case occurs, then the Jacobian of $X$ is isomorphic to the Jacobian of~$\overline X$.
\end{Lemma}

\begin{proof} We f\/irst prove that two points in the same f\/iber $p_1, p_2 \in f^{\lambda_0} = f^{-1}(\lambda_0)$ are mapped by~$T$ into a~single f\/iber~$f^{\lambda_1}$. To this end, we claim that any two points~$p_1$,~$p_2$ in dif\/ferent f\/ibers can be joined by a holomorphic section ${\mathbb C}P^1 \longrightarrow {\mathcal M}_{\rm DH}$ of the map~$f$. This can be proved quite easily by linear interpolation with sections of the form given in Lemma~\ref{sections}.

If two points $x_1$, $x_2$ of one f\/iber of $f$ are mapped by~$T$ into two dif\/ferent f\/ibers of $f$, we consider a holomorphic section~$s$ of~$f$ passing through~$T(x_1)$ and~$T(x_2)$. Then $\lambda \longmapsto T^{-1}(s(\lambda))$ is a~holomorphic immersion of~${\mathbb C}P^1$ to ${\mathcal M}_{\rm DH}$, and the degree of its normal bundle is the same as the degree of the normal bundle of the section~$s$. From Lemma~\ref{lemdeg} we obtain
\begin{gather*}(2 g_X+2)\deg (i)-2=2g_X,\end{gather*}
which yields $\deg (i)=1$. On the other hand, its degree is at least two because it passes through two points $x_1$, $x_2$ lying in one f\/iber. In view of this contradiction we conclude that~$T$ maps f\/ibers of~$f$ to f\/ibers of~$f$.

Therefore, there is a unique holomorphic map
\begin{gather*}T'\colon \ {\mathbb C}P^1 \longrightarrow {\mathbb C}P^1\end{gather*}
such that $f\circ T = T'\circ f$. As in the proof of Theorem~\ref{thm1}, $T$ must satisfy $T(\{0,\infty\}) = \{0,\infty\}$, and hence $T'$ is of the form $\lambda \longmapsto \tau\lambda$ or $\lambda \longmapsto \frac{\tau}{\lambda}$ for some $\tau \in \mathbb C^*$. If $T(0) = \infty$ the rank one Higgs moduli spaces for $X$ and $\overline X$ are holomorphically isomorphic. Since there is no nonconstant holomorphic map from an abelian variety
to an af\/f\/ine space, any biholomorphism between the rank one Higgs moduli spaces for $X$ and $\overline X$ produces a~biholomorphism between~$\operatorname{Pic}^0(X)$ and~$\operatorname{Pic}^0(\overline{X})$.
\end{proof}

Let $T \colon {\mathcal M}_{\rm DH} \longrightarrow {\mathcal M}_{\rm DH}$ be a holomorphic automorphism that maps each f\/iber of $f$ to itself, meaning $T(f^\lambda) = f^\lambda$ for all $\lambda \in {\mathbb C}P^1$. Consider the projection
\begin{gather}\label{pi2}
\pi_2 \colon \ f^0  = \operatorname{Pic}^0(X)\times H^0(X, K_X) \longrightarrow H^0(X, K_X) .
\end{gather}
Then, the automorphism $T$ restricted to $f^0$ maps f\/ibers of $\pi_2$ to f\/ibers of $\pi_2$ since $\operatorname{Pic}^0(X)$ does not have any nonconstant holomorphic map to $H^0(X, K_X)$. A~corresponding statement is true for the f\/iber $f^\infty$. Thus, the assumptions in the following
lemma are natural.

\begin{Lemma}\label{almost-id} Let $T \colon {\mathcal M}_{\rm DH} \longrightarrow {\mathcal M}_{\rm DH}$ be a~holomorphic automorphism such that
$T(f^\lambda) = f^\lambda$ for all $\lambda \in \mathbb CP^1$. Assume that $T$ restricted to $\operatorname{Pic}^0(X)\times\{0\} \subset f^0$ is the
identity map, and $T$ restricted to $\operatorname{Pic}^0(\overline{X})\times\{0\} \subset f^\infty$ is the identity map. Then, there exists an integral harmonic $1$-form $\gamma \in \Lambda = H^1(X, \mathbb Z) \subset \operatorname{Harm}(X; \mathbb C)$ such that $T$ is given by
\begin{gather*}T((\lambda, E, D)) = (\lambda, E, D+\lambda\gamma') ,\end{gather*}
where $\gamma = \gamma'+\gamma''$ is the decomposition into $(1,0)$ and $(0,1)$ components.
\end{Lemma}

\begin{proof}
Consider the ``constant'' section of $f$
\begin{gather*}s(\lambda) = [\lambda, \overline{\partial}, \lambda \partial],\end{gather*}
where $d = \partial+\overline{\partial}$ is the decomposition into types. From Lemma~\ref{sections} and the assumption that~$T$ restricted to $\operatorname{Pic}^0(X)\times\{0\} \subset f^0$ is the identity map it follows that $T\circ s$ must be of the form
\begin{gather}\label{specialT}
T\circ s(\lambda) = (\lambda, \overline{\partial}+\gamma_1'', \lambda (\partial+\gamma_2'))
\end{gather}
for some $\gamma_1, \gamma_2 \in \Lambda$. By applying the gauge $\exp\big({-}\int\gamma_1\big)$, the identity in \eqref{specialT} is equivalent to
\begin{gather*}T\circ s(\lambda) = (\lambda, \overline{\partial}, \lambda (\partial+\gamma'))\end{gather*}
 for $\gamma = \gamma_2-\gamma_1$.

Since $T$ is continuous it follows that for every constant section of~$f$ def\/ined by $s_D(\lambda) = (\lambda, L, \lambda D)$, where $D$ is a~holomorphic connection on~$L$, the equality
\begin{gather*}T\circ s_D(\lambda) = (\lambda, L, \lambda (D+\gamma'))\end{gather*}
holds. It should be mentioned that this is equivalent to
\begin{gather*}T\circ s(\lambda) = (\lambda, L\otimes L(\overline{\partial}-\gamma''), \lambda D) .\end{gather*}
Note that for any $\lambda \in {\mathbb C}\setminus\{0\}$ and for every point $p = (\lambda, E, D) \in f^\lambda$, there exists a constant section
which goes through $p$. This completes the proof.
\end{proof}

Let \begin{gather*}s \colon \ {\mathbb C}P^1 \longrightarrow {\mathcal M}_{\rm DH} ,\qquad
\lambda \longmapsto [\lambda, \overline{\partial} +\overline{\omega}+ \lambda\overline{\eta}, \lambda(\partial+\alpha)+\beta]
\end{gather*}
be a section of $f$, where $\alpha, \beta, \eta, \omega \in H^0(X, K_X)$; def\/ine
\begin{gather}\label{Ts}
T_s \colon \ {\mathcal M}_{\rm DH} \longrightarrow {\mathcal M}_{\rm DH} ,\qquad [\lambda, E, D] \longmapsto [\lambda, E\otimes
L(\overline{\partial}+\overline{\omega}+\lambda\overline{\eta}), D\otimes(\lambda(\partial+\alpha)+\beta)] ,
\end{gather}
where the tensor product of two $\lambda$-connections~$D$,~$\widetilde{D}$ on two holomorphic line bundles~$E$,~$\widetilde E$ respectively is the
$\lambda$-connection
\begin{gather*}D\otimes\widetilde{D}(e\otimes\widetilde{e}) = D(e)\otimes \widetilde{e}+e\otimes \widetilde{D}(\widetilde {e})\end{gather*}
on $E\otimes\widetilde{E}$. This $T_s$ is clearly a holomorphic automorphism of ${\mathcal M}_{\rm DH}$.

Let
\begin{gather*}
\operatorname{Aut}({\mathcal M}_{\rm DH})_0 \subset \operatorname{Aut}({\mathcal M}_{\rm DH})
\end{gather*}
be the connected component, containing the identity element, of the group of holomorphic automorphisms of ${\mathcal M}_{\rm DH}$. Using Lemmas~\ref{tau-lemma} and~\ref{almost-id} we can give the following description of $\operatorname{Aut}({\mathcal M}_{\rm DH})_0$.

\begin{Theorem}\label{thmDHa}
There is an isomorphism
\begin{gather*}\operatorname{Aut}({\mathcal M}_{\rm DH})_0 =
\mathcal M_{dR}(X)\times \big(\big(H^0(X, K_X)\times\overline{H^0(X, K_X)}\big)\rtimes {\mathbb C}^*\big)\end{gather*}
with the group structure of the right-hand side given by
\begin{gather}\label{gs37}
\big(\nabla^1, \alpha_1, \overline{\eta}_1, \tau_1\big)\cdot \big(\nabla^2, \alpha_2,\overline{\eta}_2, \tau_2\big) = \big(\nabla^1\otimes\nabla^2, \alpha_1+\tau_1\alpha_2, \overline{\eta}_1+\tfrac{1}{\tau_1}\overline{\eta}_2, \tau_1\tau_2\big) .
\end{gather}
The group $\mathcal M_{dR}(X)$ acts via the automorphisms $T_s$ in~\eqref{Ts}, $H^0(X, K_X)\times\overline{H^0(X, K_X)}$ acts using addition of forms to connections and Higgs bundles, and ${\mathbb C}^*$ acts via multiplication.
\end{Theorem}

\begin{proof} Using Lemma \ref{tau-lemma} and the lift of the $\mathbb C^*$ action on $\mathbb C P^1$ we can restrict to automorphisms $T\in \operatorname{Aut}(\mathcal M_{\rm DH})_0$ such that
\begin{gather*}T\big(f^\lambda\big) = f^\lambda\end{gather*} for all $\lambda\in\mathbb CP^1$. We claim that there exists a section $s \colon {\mathbb C}P^1 \longrightarrow {\mathcal M}_{\rm DH}$ of $f$ such that $T_s\circ T$ is of the form given in Lemma~\ref{almost-id}, where $T_s$ is constructed in~\eqref{Ts}.

Recall that $T$ acts f\/iberwise on the f\/ibers of~$\pi_2$ in~\eqref{pi2}. An automorphism of $\operatorname{Pic}^0(X)$ which is homotopic to the identity map is a translation by an element of $\operatorname{Pic}^0(X)$. Hence, there exists $\alpha, \omega \in H^0(X, K_X)$ such that
\begin{gather*}T(0, L, 0) = (0, L\otimes L(\overline{\partial}+\overline{\omega}), \alpha)\end{gather*}
for all $(0, L, 0) \in \pi_2^{-1}(0) \subset f^0$. The same argument yields that there are $\beta, \eta \in H^0(X, K_X)$ such that the automorphism $T_s$, for the section $s$ def\/ined by
\begin{gather*}s(\lambda) = (\lambda, \overline{\partial} -\overline{\omega}-\lambda\overline{\eta}, \lambda(\partial-\beta)-\alpha) ,\end{gather*}
has the desired properties.

Moreover, by applying the arguments in the proof of Lemma~\ref{almost-id}, we see that a~section~$s$ of~$f$ such that $T_s\circ T$ is of the form given in Lemma~\ref{almost-id} is unique up to tensoring with a constant section of the form $\lambda \longmapsto (\lambda, \overline{\partial}, \lambda(\partial+\gamma'))$ for some $\gamma \in H^1(X, \mathbb Z) \subset \operatorname{Harm}(X, \mathbb C)$.

Altogether, we obtain that any element in $\operatorname{Aut}({\mathcal M}_{\rm DH})_0$ is of the form $(\nabla, \alpha, \overline{\eta}, \tau)\in\mathcal M_{dR}(X)\times \big(\big(H^0(X, K_X)\times\overline{H^0(X, K_X)}\big)\rtimes {\mathbb C}^*\big)$ for the actions of the components $\mathcal M_{dR}(X)$, $H^0(X, K_X)$, $\overline{H^0(X, K_X)}$ and ${\mathbb C}^*$ on ${\mathcal M}_{\rm DH}$ as described above. The group structure of $\operatorname{Aut}({\mathcal M}_{\rm DH})_0$ is easily verif\/ied to be as in \eqref{gs37}.
\end{proof}
